\BSRAppendix{B}{用例验收场景与异常边界}
\noindent 每个 E 编号与 UC-01～UC-15 对应。场景采用“给定—当—则”结构；“本地”表示现有系统回归可核对，“目标”表示仍需第三方、BSR 业务或生产环境证据。测试通过只能证明相应路径，不能代替现状访谈。

\subsection*{B.1 核心场景}
\begin{BSRTable}
\rowcolor{TableHead}\color{white}\bfseries 编号/用例 & \color{white}\bfseries 给定、操作、可观察结果与证据状态 \\
E-01\newline UC-01 & 给定已存在的顾客邮箱和本人会话，当创建或查询档案，则重复邮箱不产生第二身份，回头客可按恒定 ID 检索，非本人不能读取私密资料。证据：本地档案与授权路径；保存期限待 BSR 决定。\\
E-02\newline UC-02 & 给定门店、起租日与归还日及已被占用的物理单件，当查询一段含重叠日期的租期，则返回整个时段可交付数量，未归还件继续占用，查询不锁库存。证据：本地日期/库存测试；09:00/16:00 仍属待批准政策。\\
E-03\newline UC-03 & 给定购物车多件装备与定位币种，当修改数量、尺码或租期，则逐项 NOK 租金、参考价值 8\% 押金和合计重新计算，旧报价失效；参考币种变化不改 NOK 合同额。证据：本地报价和浏览器购物车。\\
E-04\newline UC-04 & 给定当前损坏/超时规则，顾客未滚动阅读或未主动勾选时提交，则不得创建预约；完成同意后同账号/渠道/键同负载并发只建一单，响应丢失刷新后返回原单且不延长暂留；到期后列表、详情和重试均显示 EXPIRED、库存释放、旧报价只读，新意向重新选期。改负载冲突，缺键拒绝。现场开单同理并按员工隔离。证据：71 项本地测试及浏览器过期/丢响应路径。\\
E-05\newline UC-05 & 给定线上暂留订单，当顾客本人生成并确认当前合同，则实际发起人、门店经办人和来源分别留痕，未签不可演示付款；旧版只读且金额快照不回写。证据：本地签付回归与纸质见证；合格线上电子签署为目标。\\
E-06\newline UC-06 & 给定当前合同未签、暂留过期或经办人缺失，当顾客请求演示付款则拒绝且无流水；签当前版后并发重复提交仅产生一组租金/押金演示流水和一份备货任务。证据：本地自动与浏览器回归；第三方回调、对账和真实资金为目标。\\
E-07\newline UC-07 & 给定已确认订单、指派专员和可用单件，当逐件备货、复核、核码交付，则顾客看到相同准备进度和责任人；尺码、安全或合同守卫失败时不得借出。证据：本地准备同步、门店操作及权限测试。\\
E-08\newline UC-08 & 给定有效租赁，当在租期内升级、降级或部分归还，则新旧合同版本、逐件时间和剩余租金/押金差额留痕，重复请求不再收退。证据：本地变更测试；真实清算仍为目标。\\
E-09\newline UC-09 & 给定多件订单，接收店员工凭完整码查在租未还件，则只见实物和归还截止；提交时再次核码，接收门店固定为本人所在店。不同日期或门店逐件归还只记一次实还时间，当前店变化而归属店不被误改，未还件仍占库存；错误码、伪造接收店和额外单件均拒绝且不泄露顾客费用。证据：两店回归与浏览器分批异店接收。\\
E-10\newline UC-10 & 给定疑损单件，当技师上传对应部位照片并在 24 小时内初评，经理复核或退回，则照片、等级、责任人与版本保留；超时只提醒，不自动鉴定或收费。证据：本地照片、鉴定及提醒测试；正式费率待批准。\\
E-11\newline UC-11 & 给定全部归还且费用核定，当结算首租宽限、超时和损坏费用，则总费用、押金抵扣、可退额和补差分列；争议或回执不明不标已到账。证据：本地分项/幂等测试；真实退款与补差为目标。\\
E-12\newline UC-12 & 给定故障单件，当技师完成诊断、15 分钟工时粒度和安全检查后提交，则经理通过安全验收才恢复可租；退回或报废继续隔离。证据：本地维修与验收测试。\\
E-13\newline UC-13 & 给定可调拨单件和有效门店，当经理发出并由目标店接收，则库存流水、当前店与永久归属更新可追溯；未接收仍在途不可租。证据：本地调拨测试；正式多店运营待验收。\\
E-14\newline UC-14 & 给定员工和当日班次，当本人上下班打卡、经理看排班与经营负责人查年报，则计划与事实分列，重复打卡和改已打卡班次拒绝。证据：本地排班/打卡测试；奖金口径待 BSR 确认。\\
E-15\newline UC-15 & 给定两店新单和不同岗位会话，店员/经理只见本店逐单提醒，经营/管理员仅见匿名门店状态计数，技师只见本人任务；已读仅影响本人，未解决事项仍保留。按订单、合同、鉴定和照片 ID 探测外店对象与不存在对象同类拒绝；权限和关键变更留审计。证据：两店多角色接口与浏览器测试；外部通知为目标。\\
\end{BSRTable}

\subsection*{B.2 必须单独复核的异常边界}
\begin{BSRTable}
\rowcolor{TableHead}\color{white}\bfseries 边界 & \color{white}\bfseries 预期结果与验证状态 \\
无效租期 & 空日期、倒置、超 30 天、门店关闭或夏令时切换时拒绝错误输入；不自动替顾客换日期。日期与库存本地路径已测，夏令时须专项复核。\\
竞争占用 & 两个顾客看到同一剩余量后同时提交，最多一个获得最后一件；失败者得到可调整的门店/尺码/日期。服务端事务路径已测，跨实例并发待生产测试。\\
预约请求重试 & 同账号、同渠道、同标识和同正文重试返回原单当前状态；同标识改顾客、门店或装备返回冲突。缺失或非法标识、失败请求不创建映射。已取消或过期原单不得借重放重建。浏览器刷新找回与服务端并发均已测。\\
重复交付/归还 & 取货码已用、过期、身份不符或单件已交付/已还，应拒绝重复写入并保留原事件；相关本地路径已测。\\
旧订单政策 & 旧合同缺新费率或同意版本时不得自动追收超时/损坏费用；进入人工核实。相关本地回归已测。\\
受损争议 & 没有对应部位照片、技师越权、经理未核定或顾客提出异议时，不扣押金且不自动结清。相关本地路径已测，正式争议期限待批准。\\
外部回执不明 & 第三方支付、退款或签署超时应保留处理中/待核实状态，不凭本地按钮宣称银行到账；真实接口未接入，属于目标验收。\\
迁移与恢复 & 在数据库副本运行两次升级，校验数量、总额、键和历史版本；失败时回到备份。升级幂等有本地证据，生产恢复演练待实施。\\
\end{BSRTable}
